% Part: first-order-logic
% Chapter: syntax-and-semantics
% Section: substitution

\documentclass[../../../include/open-logic-section]{subfiles}

\begin{document}

\olfileid{fol}{syn}{sub}

\olsection{代入}

\begin{defn}[项中的代入]
我们递归地定义 $\Subst{s}{t}{x}$，即将 $t$ \emph{代入} $s$ 中 $x$ 的每一处出现所得的结果：
\begin{enumerate}
\item \indcase{s}{c}{$\Subst{\indfrm}{t}{x}$ 即为 $s$。}

\item \indcase{s}{y}{若 $y$ 为变元且 $y \not\ident x$，则 $\Subst{\indfrm}{t}{x}$ 亦为~$s$。}

\item \indcase{s}{x}{$\Subst{\indfrm}{t}{x}$ 即为~$t$。}

\item\indcase{s}{\Atom{f}{t_1, \dots, t_n}}{$\Subst{\indfrmp}{t}{x}$ 即为
  $\Atom{f}{\Subst{t_1}{t}{x}, \dots, \Subst{t_n}{t}{x}}$。}
\end{enumerate}
\end{defn}

\begin{defn}
若 $x$ 在公式~$!A$ 中的所有自由出现均不在约束 $t$ 中变元的量词的辖域内，则称项~$t$ 在 $!A$ 中对 $x$ 是\emph{自由}的。
\end{defn}

\begin{ex} ~
\begin{enumerate}
\item $\Obj v_8$ 在 $\lexists[\Obj
  v_3]\Atom{\Obj A^2_4}{\Obj v_3,\Obj v_1}$ 中对 $\Obj v_1$ 自由。

\item $\Obj f^2_1(\Obj v_1, \Obj v_2)$ 在 $\lforall[\Obj v_2]\Atom{\Obj A^2_4}{\Obj v_0,\Obj v_2}$ 中对 $\Obj v_0$ \emph{不}自由。
\end{enumerate}
\end{ex}

\begin{defn}[公式中的代入]
设 $!A$ 为一公式，$x$ 为一变元，且 $t$ 是在 $!A$ 中对 $x$ 自由的项，则 $\Subst{!A}{t}{x}$ 表示将 $t$ 代入 $!A$ 中 $x$ 的所有自由出现所得的结果。
\begin{enumerate}
\tagitem{prvFalse}{\indcase{!A}{\lfalse}{$\Subst{\indfrm}{t}{x}$ 即为 $\lfalse$。}}{}

\tagitem{prvTrue}{\indcase{!A}{\ltrue}{$\Subst{\indfrm}{t}{x}$ 即为 $\ltrue$。}}{}

\item \indcase{!A}{\Atom{P}{t_1,\dots,
    t_n}}{$\Subst{\indfrm}{t}{x}$ 即为 $\Atom{P}{\Subst{t_1}{t}{x},
    \dots, \Subst{t_n}{t}{x}}$。}

\item \indcase{!A}{\eq[t_1][t_2]}{$\Subst{\indfrmp}{t}{x}$ 即为
  $\Subst{t_1}{t}{x} = \Subst{t_2}{t}{x}$。}

\tagitem{prvNot}{\indcase{!A}{\lnot !B}{$\Subst{\indfrmp}{t}{x}$ 即为
    $\lnot \Subst{!B}{t}{x}$。}}{}

\tagitem{prvAnd}{\indcase{!A}{(!B \land
    !C)}{$\Subst{\indfrmp}{t}{x}$ 即为 $(\Subst{!B}{t}{x} \land
    \Subst{!C}{t}{x})$。}}{}

\tagitem{prvOr}{\indcase{!A}{(!B \lor
    !C)}{$\Subst{\indfrmp}{t}{x}$ 即为 $(\Subst{!B}{t}{x} \lor
    \Subst{!C}{t}{x})$。}}{}

\tagitem{prvIf}{\indcase{!A}{(!B \lif
    !C)}{$\Subst{\indfrmp}{t}{x}$ 即为 $(\Subst{!B}{t}{x} \lif
    \Subst{!C}{t}{x})$。}}{}

\tagitem{prvIff}{\indcase{!A}{(!B \liff
    !C)}{$\Subst{\indfrmp}{t}{x}$ 即为 $(\Subst{!B}{t}{x} \liff
    \Subst{!C}{t}{x})$。}}{}

\tagitem{prvAll}{
\indcase{!A}{\lforall[y][!B]}{$\Subst{\indfrmp}{t}{x}$
    即为 $\lforall[y][\Subst{!B}{t}{x}]$，若 $y$ 为不同于 $x$ 的变元；
    否则 $\Subst{\indfrmp}{t}{x}$ 即为 $\indfrm$。}}{}

\tagitem{prvEx}{
\indcase{!A}{\lexists[y][!B]}{$\Subst{\indfrmp}{t}{x}$
    即为 $\lexists[y][\Subst{!B}{t}{x}]$，若 $y$ 为不同于 $x$ 的变元；
    否则 $\Subst{\indfrmp}{t}{x}$ 即为 $\indfrm$。}}{}
\end{enumerate}
\end{defn}

\begin{explain}
注意代入可能是空代入：如果 $x$ 根本不在 $!A$ 中出现，那么 $\Subst{!A}{t}{x}$ 即为~$!A$。

要求 $t$ 必须在 $!A$ 中对 $x$ 自由，这一限制是必要的，它可以排除以下情况。若 $!A \ident \lexists[y][x < y]$ 且 $t \ident y$，则 $\Subst{!A}{t}{x}$ 将为 $\lexists[y][y < y]$。此时，自由变元 $y$ 在代入后被量词 $\lexists[y]$ “捕获”了，这是不可取的。例如，我们希望若 $\lforall[x][!B]$ 成立，则 $\Subst{!B}{t}{x}$ 亦成立。但考虑 $\lforall[x][\lexists[y][x < y]]$（这里 $!B$ 是 $\lexists[y][x < y]$）。这是一个语句，例如，它在自然数中为真：对每个数 $x$，都存在大于它的数 $y$。若我们允许用 $y$ 代入 $x$，就会得到 $\Subst{!B}{y}{x} \ident \lexists[y][y < y]$，而这是假的。我们通过要求 $t$ 中的自由变元在代入 $!A$ 后不会被 $!A$ 中的量词所约束，来避免这种情况。
\end{explain}

我们常常使用下述约定来避免繁琐的记法：如果公式 $!A$ 可能含有自由变元 $x$，我们也将其写作 $!A(x)$ 以作标识。当我们明确所讨论的 $!A$ 和 $x$ 为何，且 $t$ 是一个项（假定其在 $!A(x)$ 中对 $x$ 自由）时，我们便用 $!A(t)$ 作为 $\Subst{!A}{t}{x}$ 的简写。因此，例如我们可能会说，“称 $!A(t)$ 为 $\lforall[x][!A(x)]$ 的一个\emph{实例}。” 我们的意思是，若 $!A$ 为任意公式，$x$ 为变元，且 $t$ 是在 $!A$ 中对 $x$ 自由的项，那么 $\Subst{!A}{t}{x}$ 就是 $\lforall[x][!A]$ 的一个实例。
\end{document}